\section{Method}\label{sec:method}

\subsection{Filters and worlds}

\paragraph{Localization.} The localization problem is the running example of the particle filter tutorial by Elfring
et al.\ \cite{elfring2021particle}, simulated with its own classes. The world is a $10 \times 10$\,m square with four
landmarks at $(2, 2)$, $(2, 8)$, $(9, 2)$ and $(8, 9)$\,m; positions wrap around at its edges. The robot starts at
$(7.5, 2)$\,m facing along the $y$ axis and makes 30 steps, each commanding a forward motion of 0.25\,m and a turn
of 0.02\,rad. Its true motion has noise with standard deviations 0.005\,m and 0.002\,rad, and after each step it
measures the range and bearing to all four landmarks with noise of 0.2\,m and 0.05\,rad. The filter assumes larger
noise, 0.1\,m and 0.2\,rad in motion and 0.4\,m and 0.3\,rad in measurement, the settings of the tutorial. It does
not know the start: the particles are drawn uniformly over the world and all headings, so the filter solves global
localization during the first steps.

\paragraph{SLAM.} The SLAM problem is the FastSLAM scenario of PythonRobotics \cite{sakai2018pythonrobotics},
simulated with its own functions. The robot waits for 3\,s and then drives at 1\,m/s with a yaw rate of 0.1\,rad/s for 50\,s in total, 500 steps of 0.1\,s,
among eight landmarks. It observes the range and bearing of every landmark within 20\,m, with noise of 0.3\,m and
$2^\circ$, and the filter receives the control with noise of 0.5\,m/s and $10^\circ$/s plus a constant yaw-rate
offset of 0.01\,rad/s. The filter assumes noise of 3\,m and $10^\circ$ for observations and 1\,m/s and $20^\circ$/s for
controls, the settings of the upstream script. All particles start at the true initial pose, and the landmark
identities are known. Each particle carries one extended Kalman filter per landmark it has observed.

\paragraph{Seeds.} Every configuration runs with 20 seeds. A seed fixes both the simulated world (the true path and
all measurements) and the random numbers of the filter, and every variant of an experiment uses the same seeds, so
differences between variants within a seed come from the sampling choice alone.

\subsection{Implementation}

The filters are built on two open-source implementations: the Python particle filter tutorial of Elfring et al.\
\cite{elfring2021particle} for localization and the four resampling schemes, and the FastSLAM 1.0 and 2.0 scripts of
PythonRobotics \cite{sakai2018pythonrobotics}. Both are MIT-licensed and pinned to fixed commits (\cref{app:upstream}). The four resampling algorithms
and the three rules for when to resample are called unchanged from the tutorial. The motion and measurement models,
the proposals and the FastSLAM updates are copied, vectorized over particles and changed to compute weights in log
space.

Each filter runs the same loop at every step: the proposal draws new poses and returns log weights, the weights are
normalized, the effective sample size \eqref{eq:ess} is recorded, the rule decides whether to resample, and if so the
resampling scheme selects the particles. Proposals, resampling schemes and rules are interchangeable components
selected by name, so an experiment is a configuration file that lists the settings to keep fixed and the settings to
vary. Every run is stored with its settings, metrics, per-step values and a fingerprint of the code that produced it.

Before any experiment, each copied function was compared with the original. With the corrections switched off, the
copies reproduce the original results: the localization likelihood and motion model exactly, one full extended Kalman
particle filter update, and the FastSLAM prediction and observation updates over several steps. The comparison also
exposed errors in the original code, among them a FastSLAM 2.0 update that never samples from its proposal and an
extended Kalman particle filter whose weights break at the world's edges. The experiments use the corrected versions,
described in \cref{app:upstream}.

The experiments ran in a container with Python 3.11 and NumPy 2.2, with 16 runs in parallel on a 52-core computer.
Runtimes are therefore comparable within an experiment but only indicative in absolute terms.

\subsection{Metrics}

Let $\hat{\x}_t$ be the weighted mean of the particles after each step (the heading as a circular mean) and $\x_t$ the
true pose, over $T$ steps.
\begin{itemize}
  \item \emph{Position RMSE}: $\sqrt{\frac{1}{T} \sum_t \lVert \hat{p}_t - p_t \rVert^2}$, with $p$ the position part of
        the pose.
  \item \emph{ATE} (absolute trajectory error): the same after the best rigid alignment (rotation and translation)
        of the estimated path onto the true one. In SLAM the map frame is only anchored at the start, so a rotation
        of the whole solution is not an error in its shape; the ATE removes it. In localization it is close to the
        RMSE.
  \item \emph{Map error (aligned)}: the mean distance between the estimated and true landmark positions at the end of
        a SLAM run, after the best rigid alignment of the estimated map.
  \item \emph{Mean ESS\,/\,N}: the effective sample size before resampling, divided by $N$ and averaged over steps.
  \item \emph{Distinct particles after resampling\,/\,N}: the number of different particles that survive a
        resampling step, divided by $N$ and averaged over the steps that resample. It measures impoverishment.
  \item \emph{Resampling rate}: the fraction of steps that resample.
  \item \emph{Median NEES}: the normalized estimation error squared, $e_t^\top S_t^{-1} e_t$ with $e_t$ the pose error
        and $S_t$ the weighted covariance of the particles, whose median over the steps of a run is reported. For a
        filter whose spread matches its error the NEES is about 3, the pose dimension; much larger values mean the
        particles are overconfident. The median is used because the first steps after a uniform start give values
        of $10^8$ and more that would dominate a mean.
  \item \emph{Runtime per step}: wall-clock time of the filter step, without simulation and logging.
\end{itemize}

Each table reports the mean and standard deviation over the 20 seeds. A difference between two variants is called
clear when their 95\% confidence intervals for the mean, $\bar{m} \pm 1.96\, s / \sqrt{20}$, do not overlap, and within
the seed-to-seed spread otherwise. This test is conservative, and it is applied only between variants with the same
value of the other varied setting, for example between resampling schemes at the same number of particles.
